\section{Introduction}\label{sec:intro}

Positive scalar curvature is expected to force a manifold to lose two
dimensions at a controlled geometric scale. A spectral lower bound
mixes scalar curvature with the energy of a test function and permits
negative scalar curvature on parts of the manifold. We show that the
uniform width conclusion survives this weakening. The geometric input
is the stabilized width theorem of \cite{OpenAIWidth}, stated precisely
in Theorem~\ref{interface:width}; the work here constructs its hypotheses
from the spectral inequality.

Let $(M^n,g)$ be a connected complete smooth Riemannian manifold without
boundary. We use $\Delta_g=\diver_g\nabla_g$ and integrate against the
Riemannian volume density. The \emph{normalized spectral scalar-curvature
bound} is
\begin{equation}\label{eq:spectral-bound}
 \int_M\bigl(4|\nabla_g\phi|_g^2+\Scal_g\phi^2\bigr)\,\dd\vol_g
 \ge \int_M\phi^2\,\dd\vol_g
 \qquad\bigl(\phi\in C_c^\infty(M;\R)\bigr).
\end{equation}
Thus $-4\Delta_g+\Scal_g\ge1$ in the quadratic-form sense. This is
the definition on noncompact manifolds as well; it does not require
a lowest eigenvalue to be attained.

For a metric space $X$ and an integer $d\ge0$, its \emph{Urysohn
$d$-width} is
\[
 \UW_d(X)=\inf_{f:X\to K}\ \sup_{y\in K}\diam_X f^{-1}(y),
\]
where $f$ is continuous and $K$ is a simplicial complex of dimension
at most $d$ with its usual polyhedral topology. The diameter refers
to the entire fiber, including all its components; empty fibers may
be omitted.

\begin{theorem}\label{thm:main}
For every integer $n\ge4$ there is a finite constant $C_n$ with the
following property. Every connected complete smooth Riemannian
$n$-manifold $(M,g)$ without boundary satisfying
\eqref{eq:spectral-bound} admits a continuous map $f:M\to K$ to a
simplicial complex of dimension at most $n-2$ such that
\[
 \diam_g f^{-1}(y)\le C_n\qquad(y\in K).
\]
In particular, $\UW_{n-2}(M,g)\le C_n$.
\end{theorem}

The constant depends only on the dimension. The theorem applies to
compact and noncompact manifolds, without orientability, spin, or
bounded-geometry assumptions. If the right side of
\eqref{eq:spectral-bound} is multiplied by $\lambda>0$, scaling the
metric to $\lambda g$ gives the fiber bound $C_n/\sqrt\lambda$.

\subsection{History and relation to earlier work}

Gromov's macroscopic dimension program relates positive scalar curvature
to the dimension of a space seen at large scales. The quantitative
formulation asks whether a fixed positive scalar-curvature normalization
forces a dimension-dependent bound for the fibers of a map to an
$(n-2)$-dimensional polyhedron. This formulation goes back to
\cite[\S\,2.A(c)]{Gromov1986}; see also
\cite[\S\,$2\frac12$]{Gromov1996} and
\cite{Gromov2023}. Uniform width is stronger than merely
asking for a finite fiber bound for each individual metric. The
stabilized partition theorem used here is the geometric input from
\cite{OpenAIWidth} that gives the uniform conclusion.

Several earlier results illustrate the distinction between the
geometric and topological versions of the problem. Early
three-dimensional distance-level diameter estimates of
Gromov--Lawson~\cite{GromovLawson1983} use vanishing rational first
homology, as clarified in Gromov's later account~\cite{Gromov2023}.
That account also gives a uniform Urysohn $1$-width bound for arbitrary
complete three-manifolds with scalar curvature bounded below by a
positive constant. Under such a uniform lower bound,
Liokumovich--Maximo~\cite{LiokumovichMaximo2023} and
Liokumovich--Wang~\cite{LiokumovichWang2026} construct Morse functions
with simultaneous area and diameter bounds on connected level
components, in the closed and complete noncompact cases, respectively.
In dimensions four and five,
Chodosh--Li--Liokumovich proved finite Urysohn $1$-width for universal
covers of closed positive-scalar-curvature manifolds under the
respective assumptions $\pi_2=0$ and $\pi_2=\pi_3=0$
\cite[Theorem~4]{ChodoshLiLiokumovich2023}. That conclusion is
qualitative and uses the topology of a compact quotient. The setting
of Theorem~\ref{thm:main} instead allows arbitrary complete manifolds
and asks for one constant depending only on the dimension.

The coefficient $4$ in \eqref{eq:spectral-bound} is natural in weighted
scalar curvature. On closed manifolds, Perelman's $\mathcal F$-functional and its lowest
eigenvalue formulation involve $-4\Delta+\Scal$, together with the
expression $\Scal+2\Delta t-|\dd t|^2$
\cite[\S\S1--2]{Perelman2002}. Gromov's torus-stabilized scalar
curvature makes the same connection through diagonal warped products
\cite[\S\S1.A--1.B]{Gromov2024}. Spectral curvature bounds also enter
the band and cube inequalities of Hirsch--Kazaras--Khuri--Zhang
\cite{HirschKazarasKhuriZhang2024}, whose hypotheses depend on the
dimension, topology, and coefficient of the scalar-curvature term.
Those inequalities control separations of boundary sets, rather than
the fibers of a map on an arbitrary complete manifold.

Our starting point is the partition theorem in
\cite{OpenAIWidth}, including its locally finite stabilized input.
The spectral-to-positive-solution passage and the warped-product
formula are established tools. The additional work is to create a
quadratic reserve from an unrestricted logarithmic drift and to adapt
the local stretching argument to a signed drift. These constructions
retain the original metric as the metric in which the final fibers
are measured. In particular, no control of a conformal factor or of
the weight at infinity is needed.

\subsection{The new constructions and the proof}

The spectral hypothesis supplies a positive smooth function $v$ with
$-4\Delta_g v+\Scal_gv\ge v$. For a metric $h$ and a smooth real
function $t$, write
\begin{equation}\label{eq:drift-scalar}
 \Dcal(h,t)=\Scal_h+2\Delta_h t-|\dd t|_h^2.
\end{equation}
Then $w=-2\log v$ satisfies $\Dcal(g,w)\ge1$.
This identity produces a weighted curvature bound, but gives no
control on the value or oscillation of $w$. The stretching construction
needs an additional reserve that grows quadratically with its drift.

The first new ingredient supplies that reserve without estimating $w$.
For an arbitrary complete metric $h$ and smooth $w$, we construct a
smooth function $u$, a global graph metric $h^+=h+\dd u^2$, and a new signed drift $t^+$ such
that
\[
 \Dcal(h^+,t^+)\ge\Dcal(h,w)+(t^+)^2-C(n).
\]
Here $C(n)$ is independent of the geometry and of $w$. The graph
equation has pressure proportional to the discrepancy $u-w$; subtracting
the graph height in the drift update turns that discrepancy into the
required square. Local barriers and an angle estimate give a global
solution even when $w$ and the geometry are unbounded at infinity.
The precise statement and proof are in Proposition~\ref{pre:reserve}.

The second ingredient stretches selected pairs of sets by increasing
the metric on a compact neighborhood. Its pressure depends on a
possibly negative drift. For a dimension-dependent $T$, the reserve
is $V(t)=t^2+T^2$. When the drift increases from $t$ to $t'\ge t$,
the local estimate retains the positive part of
$V(t')-V(t)$, because increasing a signed drift need not increase
$V$. Its loss is paid by a decrease of $\tr_h h_0$, where $h_0$ is
fixed during one round of stretching, together with an arbitrarily
small summable error. The inverse construction of the pressure is
given in Proposition~\ref{pressure:law}, and the resulting geometric estimate
is Lemma~\ref{stretch:local}.

To assemble these tools, choose a radius $D$ depending only on $n$
and maximal $D$-separated centers $x_i$ in the original metric.
Their closed radius-$D$ balls cover $M$, and their closed radius-$2D$
balls are locally finite. The local lemma is applied in rounds until
each inner ball is as far as prescribed from the exterior of its
outer ball. The trace decreases
telescope within each round, independently of the number of overlapping
balls. Each index $i$ requires only finitely many rounds, and every compact
set meets only finitely many outer balls, so the final metric and drift are
smooth and eventually unchanged on each compact neighborhood.

The final bound $\Dcal(h,t)\ge1/2$ is converted into the stabilized
scalar-curvature inequality by expressing $-t$ as a locally finite sum
of compactly supported logarithmic warps and subdividing each summand
until the sum of squared gradients is small. Theorem~\ref{interface:width}
then supplies a map whose entire fibers are bounded in $g$.
Section~\ref{assembly:section} carries out the parameter choice and this
last passage. Thus the graph reserve and the signed stretching estimate
are the analytic contributions; the topological dimension drop is the
conclusion of the stated geometric input.
